% Options for packages loaded elsewhere
% Options for packages loaded elsewhere
\PassOptionsToPackage{unicode}{hyperref}
\PassOptionsToPackage{hyphens}{url}
\PassOptionsToPackage{dvipsnames,svgnames,x11names}{xcolor}
%
\documentclass[
  letterpaper,
  DIV=11,
  numbers=noendperiod]{scrartcl}
\usepackage{xcolor}
\usepackage{amsmath,amssymb}
\setcounter{secnumdepth}{-\maxdimen} % remove section numbering
\usepackage{iftex}
\ifPDFTeX
  \usepackage[T1]{fontenc}
  \usepackage[utf8]{inputenc}
  \usepackage{textcomp} % provide euro and other symbols
\else % if luatex or xetex
  \usepackage{unicode-math} % this also loads fontspec
  \defaultfontfeatures{Scale=MatchLowercase}
  \defaultfontfeatures[\rmfamily]{Ligatures=TeX,Scale=1}
\fi
\usepackage{lmodern}
\ifPDFTeX\else
  % xetex/luatex font selection
  \setmainfont[]{DejaVu Serif}
  \setmonofont[]{DejaVu Sans Mono}
\fi
% Use upquote if available, for straight quotes in verbatim environments
\IfFileExists{upquote.sty}{\usepackage{upquote}}{}
\IfFileExists{microtype.sty}{% use microtype if available
  \usepackage[]{microtype}
  \UseMicrotypeSet[protrusion]{basicmath} % disable protrusion for tt fonts
}{}
\makeatletter
\@ifundefined{KOMAClassName}{% if non-KOMA class
  \IfFileExists{parskip.sty}{%
    \usepackage{parskip}
  }{% else
    \setlength{\parindent}{0pt}
    \setlength{\parskip}{6pt plus 2pt minus 1pt}}
}{% if KOMA class
  \KOMAoptions{parskip=half}}
\makeatother
% Make \paragraph and \subparagraph free-standing
\makeatletter
\ifx\paragraph\undefined\else
  \let\oldparagraph\paragraph
  \renewcommand{\paragraph}{
    \@ifstar
      \xxxParagraphStar
      \xxxParagraphNoStar
  }
  \newcommand{\xxxParagraphStar}[1]{\oldparagraph*{#1}\mbox{}}
  \newcommand{\xxxParagraphNoStar}[1]{\oldparagraph{#1}\mbox{}}
\fi
\ifx\subparagraph\undefined\else
  \let\oldsubparagraph\subparagraph
  \renewcommand{\subparagraph}{
    \@ifstar
      \xxxSubParagraphStar
      \xxxSubParagraphNoStar
  }
  \newcommand{\xxxSubParagraphStar}[1]{\oldsubparagraph*{#1}\mbox{}}
  \newcommand{\xxxSubParagraphNoStar}[1]{\oldsubparagraph{#1}\mbox{}}
\fi
\makeatother

\usepackage{color}
\usepackage{fancyvrb}
\newcommand{\VerbBar}{|}
\newcommand{\VERB}{\Verb[commandchars=\\\{\}]}
\DefineVerbatimEnvironment{Highlighting}{Verbatim}{commandchars=\\\{\}}
% Add ',fontsize=\small' for more characters per line
\usepackage{framed}
\definecolor{shadecolor}{RGB}{241,243,245}
\newenvironment{Shaded}{\begin{snugshade}}{\end{snugshade}}
\newcommand{\AlertTok}[1]{\textcolor[rgb]{0.68,0.00,0.00}{#1}}
\newcommand{\AnnotationTok}[1]{\textcolor[rgb]{0.37,0.37,0.37}{#1}}
\newcommand{\AttributeTok}[1]{\textcolor[rgb]{0.40,0.45,0.13}{#1}}
\newcommand{\BaseNTok}[1]{\textcolor[rgb]{0.68,0.00,0.00}{#1}}
\newcommand{\BuiltInTok}[1]{\textcolor[rgb]{0.00,0.23,0.31}{#1}}
\newcommand{\CharTok}[1]{\textcolor[rgb]{0.13,0.47,0.30}{#1}}
\newcommand{\CommentTok}[1]{\textcolor[rgb]{0.37,0.37,0.37}{#1}}
\newcommand{\CommentVarTok}[1]{\textcolor[rgb]{0.37,0.37,0.37}{\textit{#1}}}
\newcommand{\ConstantTok}[1]{\textcolor[rgb]{0.56,0.35,0.01}{#1}}
\newcommand{\ControlFlowTok}[1]{\textcolor[rgb]{0.00,0.23,0.31}{\textbf{#1}}}
\newcommand{\DataTypeTok}[1]{\textcolor[rgb]{0.68,0.00,0.00}{#1}}
\newcommand{\DecValTok}[1]{\textcolor[rgb]{0.68,0.00,0.00}{#1}}
\newcommand{\DocumentationTok}[1]{\textcolor[rgb]{0.37,0.37,0.37}{\textit{#1}}}
\newcommand{\ErrorTok}[1]{\textcolor[rgb]{0.68,0.00,0.00}{#1}}
\newcommand{\ExtensionTok}[1]{\textcolor[rgb]{0.00,0.23,0.31}{#1}}
\newcommand{\FloatTok}[1]{\textcolor[rgb]{0.68,0.00,0.00}{#1}}
\newcommand{\FunctionTok}[1]{\textcolor[rgb]{0.28,0.35,0.67}{#1}}
\newcommand{\ImportTok}[1]{\textcolor[rgb]{0.00,0.46,0.62}{#1}}
\newcommand{\InformationTok}[1]{\textcolor[rgb]{0.37,0.37,0.37}{#1}}
\newcommand{\KeywordTok}[1]{\textcolor[rgb]{0.00,0.23,0.31}{\textbf{#1}}}
\newcommand{\NormalTok}[1]{\textcolor[rgb]{0.00,0.23,0.31}{#1}}
\newcommand{\OperatorTok}[1]{\textcolor[rgb]{0.37,0.37,0.37}{#1}}
\newcommand{\OtherTok}[1]{\textcolor[rgb]{0.00,0.23,0.31}{#1}}
\newcommand{\PreprocessorTok}[1]{\textcolor[rgb]{0.68,0.00,0.00}{#1}}
\newcommand{\RegionMarkerTok}[1]{\textcolor[rgb]{0.00,0.23,0.31}{#1}}
\newcommand{\SpecialCharTok}[1]{\textcolor[rgb]{0.37,0.37,0.37}{#1}}
\newcommand{\SpecialStringTok}[1]{\textcolor[rgb]{0.13,0.47,0.30}{#1}}
\newcommand{\StringTok}[1]{\textcolor[rgb]{0.13,0.47,0.30}{#1}}
\newcommand{\VariableTok}[1]{\textcolor[rgb]{0.07,0.07,0.07}{#1}}
\newcommand{\VerbatimStringTok}[1]{\textcolor[rgb]{0.13,0.47,0.30}{#1}}
\newcommand{\WarningTok}[1]{\textcolor[rgb]{0.37,0.37,0.37}{\textit{#1}}}

\usepackage{longtable,booktabs,array}
\usepackage{calc} % for calculating minipage widths
% Correct order of tables after \paragraph or \subparagraph
\usepackage{etoolbox}
\makeatletter
\patchcmd\longtable{\par}{\if@noskipsec\mbox{}\fi\par}{}{}
\makeatother
% Allow footnotes in longtable head/foot
\IfFileExists{footnotehyper.sty}{\usepackage{footnotehyper}}{\usepackage{footnote}}
\makesavenoteenv{longtable}
\usepackage{graphicx}
\makeatletter
\newsavebox\pandoc@box
\newcommand*\pandocbounded[1]{% scales image to fit in text height/width
  \sbox\pandoc@box{#1}%
  \Gscale@div\@tempa{\textheight}{\dimexpr\ht\pandoc@box+\dp\pandoc@box\relax}%
  \Gscale@div\@tempb{\linewidth}{\wd\pandoc@box}%
  \ifdim\@tempb\p@<\@tempa\p@\let\@tempa\@tempb\fi% select the smaller of both
  \ifdim\@tempa\p@<\p@\scalebox{\@tempa}{\usebox\pandoc@box}%
  \else\usebox{\pandoc@box}%
  \fi%
}
% Set default figure placement to htbp
\def\fps@figure{htbp}
\makeatother


% definitions for citeproc citations
\NewDocumentCommand\citeproctext{}{}
\NewDocumentCommand\citeproc{mm}{%
  \begingroup\def\citeproctext{#2}\cite{#1}\endgroup}
\makeatletter
 % allow citations to break across lines
 \let\@cite@ofmt\@firstofone
 % avoid brackets around text for \cite:
 \def\@biblabel#1{}
 \def\@cite#1#2{{#1\if@tempswa , #2\fi}}
\makeatother
\newlength{\cslhangindent}
\setlength{\cslhangindent}{1.5em}
\newlength{\csllabelwidth}
\setlength{\csllabelwidth}{3em}
\newenvironment{CSLReferences}[2] % #1 hanging-indent, #2 entry-spacing
 {\begin{list}{}{%
  \setlength{\itemindent}{0pt}
  \setlength{\leftmargin}{0pt}
  \setlength{\parsep}{0pt}
  % turn on hanging indent if param 1 is 1
  \ifodd #1
   \setlength{\leftmargin}{\cslhangindent}
   \setlength{\itemindent}{-1\cslhangindent}
  \fi
  % set entry spacing
  \setlength{\itemsep}{#2\baselineskip}}}
 {\end{list}}
\usepackage{calc}
\newcommand{\CSLBlock}[1]{\hfill\break\parbox[t]{\linewidth}{\strut\ignorespaces#1\strut}}
\newcommand{\CSLLeftMargin}[1]{\parbox[t]{\csllabelwidth}{\strut#1\strut}}
\newcommand{\CSLRightInline}[1]{\parbox[t]{\linewidth - \csllabelwidth}{\strut#1\strut}}
\newcommand{\CSLIndent}[1]{\hspace{\cslhangindent}#1}



\setlength{\emergencystretch}{3em} % prevent overfull lines

\providecommand{\tightlist}{%
  \setlength{\itemsep}{0pt}\setlength{\parskip}{0pt}}



 


\KOMAoption{captions}{tableheading}
\makeatletter
\@ifpackageloaded{caption}{}{\usepackage{caption}}
\AtBeginDocument{%
\ifdefined\contentsname
  \renewcommand*\contentsname{Table of contents}
\else
  \newcommand\contentsname{Table of contents}
\fi
\ifdefined\listfigurename
  \renewcommand*\listfigurename{List of Figures}
\else
  \newcommand\listfigurename{List of Figures}
\fi
\ifdefined\listtablename
  \renewcommand*\listtablename{List of Tables}
\else
  \newcommand\listtablename{List of Tables}
\fi
\ifdefined\figurename
  \renewcommand*\figurename{Figure}
\else
  \newcommand\figurename{Figure}
\fi
\ifdefined\tablename
  \renewcommand*\tablename{Table}
\else
  \newcommand\tablename{Table}
\fi
}
\@ifpackageloaded{float}{}{\usepackage{float}}
\floatstyle{ruled}
\@ifundefined{c@chapter}{\newfloat{codelisting}{h}{lop}}{\newfloat{codelisting}{h}{lop}[chapter]}
\floatname{codelisting}{Listing}
\newcommand*\listoflistings{\listof{codelisting}{List of Listings}}
\makeatother
\makeatletter
\makeatother
\makeatletter
\@ifpackageloaded{caption}{}{\usepackage{caption}}
\@ifpackageloaded{subcaption}{}{\usepackage{subcaption}}
\makeatother
\usepackage{fixfoot}
\usepackage{perpage}
\MakePerPage{footnote}
\long\def\rfnbodya{See Silva, José Afonso da. 2026. \emph{Curso de Direito Constitucional
Positivo}. 47th ed. São Paulo: JusPodivm.}
\DeclareFixedFootnote{\rfnnotea}{\protect\rfnbodya}
\long\def\rfnbodyb{See Brasil. 1988. {``Constituição Da República Federativa Do Brasil de
1988.''} Brasília, DF: Presidência da República.
\url{https://www.planalto.gov.br/ccivil_03/constituicao/constituicao.htm}.}
\DeclareFixedFootnote{\rfnnoteb}{\protect\rfnbodyb}
\long\def\rfnbodyc{See Bloch, Marc. 1953. \emph{The Historian's Craft}. Manchester:
Manchester University Press.}
\DeclareFixedFootnote{\rfnnotec}{\protect\rfnbodyc}
\usepackage{bookmark}
\IfFileExists{xurl.sty}{\usepackage{xurl}}{} % add URL line breaks if available
\urlstyle{same}
\hypersetup{
  pdftitle={reusable-footnotes},
  pdfauthor={Lucas Sá Barreto Jordão},
  colorlinks=true,
  linkcolor={blue},
  filecolor={Maroon},
  citecolor={Blue},
  urlcolor={Blue},
  pdfcreator={LaTeX via pandoc}}


\title{reusable-footnotes}
\usepackage{etoolbox}
\makeatletter
\providecommand{\subtitle}[1]{% add subtitle to \maketitle
  \apptocmd{\@title}{\par {\large #1 \par}}{}{}
}
\makeatother
\subtitle{Page-local reusable bibliographic footnotes for Quarto}
\author{Lucas Sá Barreto Jordão}
\date{2026-09-22}
\begin{document}
\maketitle

\renewcommand*\contentsname{Table of contents}
{
\hypersetup{linkcolor=}
\setcounter{tocdepth}{3}
\tableofcontents
}

\section{Why this pattern matters}\label{why-this-pattern-matters}

\texttt{reusable-footnotes} is designed for a citation pattern that is
especially common in \textbf{Law}. Legal scholarship, briefs, opinions,
memoranda, and books frequently place bibliographic references in page
footnotes rather than relying only on a final reference list. The same
practice is also familiar in \textbf{History, Philosophy, Theology,
Classics, literary studies, and other Humanities disciplines}.

The important distinction is between the \textbf{bibliographic source}
and its \textbf{footnote occurrences}. BibTeX should define the source
once; Quarto should decide how that source is rendered wherever it is
used.

For printable outputs, this has an important practical consequence. In
\textbf{PDF} and \textbf{DOCX}, the reader can encounter the full source
at the bottom of the same page as the proposition it supports. A
reference cited again on another page can therefore appear again on that
later page, while repeated uses on the \emph{same} page can share one
footnote.

The intended page-local behavior is:

\begin{Shaded}
\begin{Highlighting}[]
\NormalTok{page 1     A¹ ... A¹ ... B²}
\NormalTok{           ─────────────────}
\NormalTok{           ¹ Full reference A}
\NormalTok{           ² Full reference B}

\NormalTok{page 2     A³ ... A³ ... B⁴       \# continuous numbering}
\NormalTok{           ─────────────────}
\NormalTok{           ³ Full reference A}
\NormalTok{           ⁴ Full reference B}
\end{Highlighting}
\end{Shaded}

or, when numbering restarts on every page:

\begin{Shaded}
\begin{Highlighting}[]
\NormalTok{page 1     A¹ ... A¹ ... B²}
\NormalTok{page 2     A¹ ... A¹ ... B²}
\end{Highlighting}
\end{Shaded}

In both cases, the source is repeated \textbf{between pages}, but not
unnecessarily duplicated \textbf{within the same page}.

\section{BibTeX and full bibliography
entries}\label{bibtex-and-full-bibliography-entries}

A normal Pandoc/Quarto citation remains completely unchanged:

\begin{Shaded}
\begin{Highlighting}[]
\NormalTok{A normal citation stays normal: [@silva2026].}
\end{Highlighting}
\end{Shaded}

For example: a normal citation stays normal (Silva 2026).

The development version of Fred Guth's \texttt{bibentry} extension adds
a second representation:

\begin{Shaded}
\begin{Highlighting}[]
\NormalTok{[@silva2026]\{.bibentry\}}
\end{Highlighting}
\end{Shaded}

That form prints the \textbf{full bibliography entry using the
bibliography layout of the active CSL style}. It does not redefine what
\texttt{@silva2026} means.

Here is the full entry inline:

Silva, José Afonso da. 2026. \emph{Curso de Direito Constitucional
Positivo}. 47th ed. São Paulo: JusPodivm.

With:

\begin{Shaded}
\begin{Highlighting}[]
\FunctionTok{bibliography}\KeywordTok{:}\AttributeTok{ references.bib}
\FunctionTok{csl}\KeywordTok{:}\AttributeTok{ your{-}style.csl}
\end{Highlighting}
\end{Shaded}

ordinary citations use the CSL citation layout, while \texttt{.bibentry}
extracts the corresponding bibliography-layout entry.

\section{Legal-style footnotes}\label{legal-style-footnotes}

The useful composition remains native Pandoc syntax:

\begin{Shaded}
\begin{Highlighting}[]
\NormalTok{A legal proposition.\^{}[See [@silva2026]\{.bibentry\}]}
\end{Highlighting}
\end{Shaded}

Rendered here:

A constitutional proposition may be supported in a page
footnote.\rfnnotea

The Constitution itself can be rendered in full the same way.\rfnnoteb

The mechanism is not limited to Law. A historical argument can use the
same pattern.\rfnnotec

\section{Multi-page print
demonstration}\label{multi-page-print-demonstration}

The following example is intentionally spread over several physical
pages using Quarto's native \texttt{pagebreak} shortcode. This makes the
page-local behavior visible in \textbf{PDF} and \textbf{DOCX} rather
than keeping all examples on one page.

\subsection{Demonstration page 1}\label{demonstration-page-1}

A printable legal document often develops an argument over several
paragraphs while keeping authorities immediately available to the
reader. The first occurrence below creates the full footnote for José
Afonso da Silva.\rfnnotea

A later sentence on the same page can rely on exactly the same
bibliographic footnote without printing the long entry a second
time.\rfnnotea

The Constitution is a distinct source and therefore receives its own
footnote.\rfnnoteb

This page-local model is useful when the body text is intended to be
read on paper or as a paginated digital document: the source remains
close to the proposition, but repeated use of the same source does not
make the bottom of the page unnecessarily repetitive.

A final sentence on this first page cites the Constitution again.
Because the footnote content is identical and the occurrence is on the
same page, it reuses the existing note.\rfnnoteb

\newpage{}

\subsection{Demonstration page 2}\label{demonstration-page-2}

This second page deliberately cites the same authorities again. The
Silva reference must now be printed again at the bottom of this page,
because a reader should not need to return to the previous page to
identify the source.\rfnnotea

A second occurrence on this same page should reuse the new page-local
Silva footnote.\rfnnotea

The Constitution likewise receives a new page-local occurrence
here.\rfnnoteb

The point is not to treat the bibliographic source as a new source.
BibTeX still contains only one \texttt{silva2026} entry and one
\texttt{brasil1988} entry. What changes is the \textbf{placement of the
footnote instance} in the printable document.

This page ends by citing the Constitution once more, demonstrating reuse
within the second page.\rfnnoteb

\newpage{}

\subsection{Demonstration page 3}\label{demonstration-page-3}

A third page repeats the same pattern. José Afonso da Silva appears
again as a footnote on this page.\rfnnotea

The same source can be invoked again without duplicating the footnote
text on this page.\rfnnotea

A historical source can participate in exactly the same
mechanism.\rfnnotec

A second reference to the historical source on the same page reuses its
page-local note.\rfnnotec

The result is a document in which each printed page carries the
references needed to understand that page, while identical notes are not
redundantly printed twice on the same page.

\section{Configuration}\label{configuration}

The repository example uses page-local reuse and restarts footnote
numbering on every page:

\begin{Shaded}
\begin{Highlighting}[]
\FunctionTok{bibliography}\KeywordTok{:}\AttributeTok{ references.bib}

\FunctionTok{filters}\KeywordTok{:}
\AttributeTok{  }\KeywordTok{{-}}\AttributeTok{ bibentry}
\AttributeTok{  }\KeywordTok{{-}}\AttributeTok{ reusable{-}footnotes}

\FunctionTok{reusable{-}footnotes}\KeywordTok{:}
\AttributeTok{  }\FunctionTok{enabled}\KeywordTok{:}\AttributeTok{ }\CharTok{true}
\AttributeTok{  }\FunctionTok{backlinks}\KeywordTok{:}\AttributeTok{ }\CharTok{true}
\AttributeTok{  }\FunctionTok{scope}\KeywordTok{:}\AttributeTok{ page}
\AttributeTok{  }\FunctionTok{numbering}\KeywordTok{:}\AttributeTok{ page}
\AttributeTok{  }\FunctionTok{docx{-}scope}\KeywordTok{:}\AttributeTok{ pagebreak}

\FunctionTok{format}\KeywordTok{:}
\AttributeTok{  }\FunctionTok{html}\KeywordTok{:}
\AttributeTok{    }\FunctionTok{toc}\KeywordTok{:}\AttributeTok{ }\CharTok{true}
\AttributeTok{  }\FunctionTok{pdf}\KeywordTok{:}
\AttributeTok{    }\FunctionTok{toc}\KeywordTok{:}\AttributeTok{ }\CharTok{true}
\AttributeTok{    }\FunctionTok{pdf{-}engine}\KeywordTok{:}\AttributeTok{ lualatex}
\AttributeTok{  }\FunctionTok{docx}\KeywordTok{:}
\AttributeTok{    }\FunctionTok{toc}\KeywordTok{:}\AttributeTok{ }\CharTok{false}
\AttributeTok{    }\FunctionTok{reference{-}doc}\KeywordTok{:}\AttributeTok{ \_extensions/reusable{-}footnotes/reference{-}page.docx}
\end{Highlighting}
\end{Shaded}

The main options are:

\begin{itemize}
\tightlist
\item
  \texttt{scope:\ page}: identical notes are reused only within the
  current page scope and may appear again later.
\item
  \texttt{numbering:\ page}: numbering restarts at \texttt{1} on every
  printed page.
\item
  \texttt{numbering:\ continuous}: numbering continues across pages
  while reuse remains page-local.
\item
  \texttt{docx-scope:\ pagebreak}: explicit \texttt{} boundaries define
  deterministic DOCX page scopes before Word performs final layout.
\item
  \texttt{backlinks:\ true}: HTML adds discreet return links for
  repeated calls to the same local footnote.
\end{itemize}

For continuous numbering, change only:

\begin{Shaded}
\begin{Highlighting}[]
\FunctionTok{reusable{-}footnotes}\KeywordTok{:}
\AttributeTok{  }\FunctionTok{scope}\KeywordTok{:}\AttributeTok{ page}
\AttributeTok{  }\FunctionTok{numbering}\KeywordTok{:}\AttributeTok{ continuous}
\AttributeTok{  }\FunctionTok{docx{-}scope}\KeywordTok{:}\AttributeTok{ pagebreak}
\end{Highlighting}
\end{Shaded}

For DOCX page-local numbering, the supplied \texttt{reference-page.docx}
configures Word's native footnote counter to restart on each page. For
continuous DOCX numbering, omit that reference document or use a custom
continuous-numbering reference DOCX.

\section{HTML websites and books}\label{html-websites-and-books}

A standalone HTML document is one scrolling document, so there is no
physical-page concept comparable to PDF or DOCX. The whole file is
therefore one logical reusable-footnote scope.

A Quarto \textbf{website} or \textbf{HTML book} is different: each
\texttt{.qmd} page or chapter is rendered as its own HTML page. The
reusable-footnote scope starts fresh for every page, but numbering can
either restart or continue globally.

\subsection{Page-local numbering}\label{page-local-numbering}

Each rendered HTML page starts again at \texttt{1}; repeated calls
within the page reuse the local canonical footnote.

\begin{itemize}
\tightlist
\item
  \href{https://github.com/lsbjordao/quarto-reusable-footnotes/tree/main/examples/website}{Website
  example --- page-local numbering}
\item
  \href{https://github.com/lsbjordao/quarto-reusable-footnotes/tree/main/examples/book}{Book
  example --- page-local numbering}
\end{itemize}

\begin{Shaded}
\begin{Highlighting}[]
\ExtensionTok{quarto}\NormalTok{ preview examples/website}
\ExtensionTok{quarto}\NormalTok{ preview examples/book}
\end{Highlighting}
\end{Shaded}

\subsection{Global numbering}\label{global-numbering}

The footnote \textbf{scope remains page-local}, but the displayed
numbers continue across pages. In the examples below, the first page
contributes \texttt{1,\ 2}, the second \texttt{3,\ 4}, and the third
\texttt{5,\ 6}.

\begin{itemize}
\tightlist
\item
  \href{https://github.com/lsbjordao/quarto-reusable-footnotes/tree/main/examples/website-global}{Website
  example --- global numbering}
\item
  \href{https://github.com/lsbjordao/quarto-reusable-footnotes/tree/main/examples/book-global}{Book
  example --- global numbering}
\end{itemize}

\begin{Shaded}
\begin{Highlighting}[]
\ExtensionTok{quarto}\NormalTok{ preview examples/website{-}global}
\ExtensionTok{quarto}\NormalTok{ preview examples/book{-}global}
\end{Highlighting}
\end{Shaded}

Because Quarto renders each HTML page independently, continuous
project-wide numbering needs an explicit page order:

\begin{Shaded}
\begin{Highlighting}[]
\FunctionTok{reusable{-}footnotes}\KeywordTok{:}
\AttributeTok{  }\FunctionTok{scope}\KeywordTok{:}\AttributeTok{ page}
\AttributeTok{  }\FunctionTok{numbering}\KeywordTok{:}\AttributeTok{ continuous}
\AttributeTok{  }\FunctionTok{html{-}order}\KeywordTok{:}
\AttributeTok{    }\KeywordTok{{-}}\AttributeTok{ index.qmd}
\AttributeTok{    }\KeywordTok{{-}}\AttributeTok{ doctrine.qmd}
\AttributeTok{    }\KeywordTok{{-}}\AttributeTok{ constitution.qmd}
\end{Highlighting}
\end{Shaded}

\texttt{html-order} is used only to calculate the numbering offset
contributed by earlier pages. It does \textbf{not} merge their footnote
scopes. A source cited on two pages can therefore appear in both pages'
footnote areas while receiving different numbers in a continuous
sequence.

This gives four distinct, demonstrable HTML modes in the repository:
website/local, website/global, book/local, and book/global.

\section{Why BibTeX remains the source of
truth}\label{why-bibtex-remains-the-source-of-truth}

The author should not copy and paste formatted references into
\texttt{\^{}{[}...{]}}. The \texttt{.bib} entry provides the stable
bibliographic identity:

\begin{Shaded}
\begin{Highlighting}[]
\VariableTok{@book}\NormalTok{\{}\OtherTok{silva2026}\NormalTok{,}
  \DataTypeTok{author}\NormalTok{    = \{Silva, José Afonso da\},}
  \DataTypeTok{title}\NormalTok{     = \{Curso de direito constitucional positivo\},}
  \DataTypeTok{edition}\NormalTok{   = \{47\},}
  \DataTypeTok{address}\NormalTok{   = \{São Paulo\},}
  \DataTypeTok{publisher}\NormalTok{ = \{JusPodivm\},}
  \DataTypeTok{year}\NormalTok{      = \{2026\},}
  \DataTypeTok{note}\NormalTok{      = \{Coedição Malheiros\}}
\NormalTok{\}}
\end{Highlighting}
\end{Shaded}

The same entry can then be represented in three different ways without
duplicating data:

\begin{Shaded}
\begin{Highlighting}[]
\NormalTok{Normal citation: [@silva2026]}

\NormalTok{Full entry: [@silva2026]\{.bibentry\}}

\NormalTok{Full entry in a footnote: \^{}[See [@silva2026]\{.bibentry\}]}
\end{Highlighting}
\end{Shaded}

\section{\texorpdfstring{\texttt{bibentry}
integration}{bibentry integration}}\label{bibentry-integration}

This repository currently vendors an \textbf{extended development
version} of Fred Guth's
\href{https://github.com/fredguth/bibentry}{\texttt{bibentry}} extension
under \texttt{\_extensions/bibentry/} for integration testing. The
original project is MIT licensed and remains credited to Frederico Guth.

The proposed upstream changes add recursive AST support, including
\texttt{.bibentry} inside native footnotes; true inline replacement; a
single citeproc pass preserving numeric CSL ordering; and portable AST
output for HTML, PDF/LaTeX, DOCX, and Typst.

\texttt{bibentry} and \texttt{reusable-footnotes} have separate
responsibilities:

\begin{itemize}
\tightlist
\item
  \textbf{\texttt{bibentry}} resolves a citation key and renders its
  full CSL bibliography form;
\item
  \textbf{\texttt{reusable-footnotes}} decides when an identical
  footnote occurrence should be reused locally or printed again on a
  later page.
\end{itemize}

\section{Output formats}\label{output-formats}

\begin{longtable}[]{@{}
  >{\raggedright\arraybackslash}p{(\linewidth - 2\tabcolsep) * \real{0.5000}}
  >{\raggedright\arraybackslash}p{(\linewidth - 2\tabcolsep) * \real{0.5000}}@{}}
\toprule\noalign{}
\begin{minipage}[b]{\linewidth}\raggedright
Format
\end{minipage} & \begin{minipage}[b]{\linewidth}\raggedright
Page-local behavior
\end{minipage} \\
\midrule\noalign{}
\endhead
\bottomrule\noalign{}
\endlastfoot
HTML document & the whole HTML file is one logical page \\
HTML website/book & each \texttt{.qmd} page owns its footnotes;
numbering may restart or continue globally \\
PDF & true physical-page reuse; identical notes print once per page and
may reappear on later pages \\
DOCX & exact reuse within explicit pagebreak-delimited pages; Word can
restart numbering on every physical page through the supplied reference
document \\
\end{longtable}

No Python runtime or post-render Python step is required.

\section*{References}\label{references}
\addcontentsline{toc}{section}{References}

\phantomsection\label{refs}
\begin{CSLReferences}{1}{0}
\bibitem[\citeproctext]{ref-bloch1953}
Bloch, Marc. 1953. \emph{The Historian's Craft}. Manchester: Manchester
University Press.

\bibitem[\citeproctext]{ref-brasil1988}
Brasil. 1988. {``Constituição Da República Federativa Do Brasil de
1988.''} Brasília, DF: Presidência da República.
\url{https://www.planalto.gov.br/ccivil_03/constituicao/constituicao.htm}.

\bibitem[\citeproctext]{ref-silva2026}
Silva, José Afonso da. 2026. \emph{Curso de Direito Constitucional
Positivo}. 47th ed. São Paulo: JusPodivm.

\end{CSLReferences}




\end{document}
